
\section{Partial Truth Tables}

% \subsection{Partial Truth Tables} \cite[][Chapter 15, pp. 102-8]{Mag25}

Last week, we saw how we can save time by only filling in the parts of our
truth tables that matter to the evaluation of the overall formula (i.e.
to its \emph{main operator}). But these are still \emph{complete} truth tables,
it's just that many of the cells are silent because they do not matter or are
obvious.

We can simplify our approach even more if we take a step back and consider what
we are trying to do in building truth tables in the first place. They are not
interesting in their own right. But they do give us a clear method for being
thorough in our exploration of the various logical possibilities, gven the rules
of the system.

One more step back. Why do we want to explore all of the logical possibilities
anyway? We want to be thorough in this way because we have defined a number of
our logical notions (e.g. \emph{tautology}, \emph{validity}, \emph{entailment})
in terms of these possibilities. When we explore these properties, we want to
show that the sentences being evaluated are true (or false) in all of the
possibilities. In our parlance from last week, we are looking for \emph{bad
lines} in the table.

But what about showing that a sentence \emph{does not} have these properties?
Surely it will suffice to just find one bad line, a \emph{counter-example} so
to speak. And indeed it will. We can describe the landscape with this table
modified from \cite{Mag25}[Chapter 15].

\begin{center}
  \begin{tabular}{c c c c}
                  & \bf{yes} & \bf{no} & should be (always)     \\
    \hline
    tautology     &  all     &  one    & true                   \\
    contradiction &  all     &  one    & false                  \\
    \hline
    equivalent    &  all     &  one    & same                   \\
    unsatisfiable &  all     &  one    & false                  \\
    \hline
    valid         &  all     &  one    & true conclusion or false premise  \\
    entailment    &  all     &  one    & true conclusion or false premise  \\
  \end{tabular}
\end{center}

Recall from \ref{def:tautology} that a TFL sentence is a tautology just in
case it is true in all scenarios. Showing that a TFL sentence is not a
tautology amounts to finding a scenario where it is not true. Since we know
we are trying to find a bad (i.e. false) line, we work backward through the
columns with this assumption in mind.

\begin{center}
  \begin{tabular}{c c|c c c c c c c}
    A & B & (A & \eif   & B) & \eif   & (B & \eif & A)  \\
    \hline
    F & T &    & T      &    & \bf{F} &  T & F    & F   \\
    4 & 4 &    & 2      &    &   1    &  3 & 2    & 3   \\
  \end{tabular}
\end{center}

The TFL sentence $(A \eif B) \eif (B \eif A)$ is false when $A$ is false and
$B$ is true. It is not a tautology.

Similarly, we can show that a TFL sentence is not a contradiction by showing
that it is sometimes true. So this time we start by assuming a true line.

\begin{center}
  \begin{tabular}{c c|c c c c c c c}
    A & B & (A & \eif   & B) & \eif   & (B & \eif & A)  \\
    \hline
    T & F &  T & F      & F  & \bf{T} &    &      &     \\
    4 & 4 &  3 & 2      & 3  &   1    &    &      &     \\
  \end{tabular}
\end{center}

This time we had to decide what to do next, since there are two ways to make
the main operator true: let its antecedant be false or let its consequent be
true. We chose the former, because that choice (a false material implication)
forces our truth value assignment to the atomic letters: $A$ must be true and
$B$ must be false. We have found a counter-example. The TFL sentence $(A \eif
B) \eif (B \eif A)$ is also not a contradiction.

We can extend this practice to truth tables with multiple TFL sentences in
order to find counter-examples for equivalence, satisfiability, validity and
entailment. For example, the following partial table demonstrates that $A \eand
B$ and $A \eand \enot B$ are not equivalent.

\begin{center}
  \begin{tabular}{c c||c c c|c c c c}
    A & B    &    A & \eand  & B    &    A & \eand & \enot & B \\
    \hline
    T & T    &      & \bf{T} &      &      & \bf{F} &  F    &   \\
    2 & 2    &      &   1    &      &      &  4     &  3    &   \\
  \end{tabular}
\end{center}

When we set the values of $A$ and $B$ to make $A \eand B$ true, the truth value
of $A \eand \enot B$ correspondingly becomes false. The two TFL sentences have
different truth values in some scenarios. They cannot be equivalent. A similar
process can be applied to the remaining logical properties.

TODO: satisfiability, validity and entailment tables ??

